\chapter{JWT Authentication}

Once a user's password is verified, the server needs a way to remember that
this browser is "logged in" on every subsequent request -- without querying
the database on every single request just to check a session table. A
\textbf{JSON Web Token (JWT)} solves this by packing identity information
into a signed, self-contained string that the server can verify without
storing anything.

\section{Anatomy of a JWT}

A JWT is three base64url-encoded segments separated by dots:
\texttt{header.payload.signature}.

\begin{center}
\begin{tabularx}{\textwidth}{lX}
\toprule
\textbf{Segment} & \textbf{Contents} \\
\midrule
Header & Metadata: the signing algorithm (e.g. \texttt{HS256}) and token type. \\
Payload & Claims -- the actual data, e.g. \texttt{\{ id: userId, iat, exp \}}. Readable by anyone, not encrypted. \\
Signature & \texttt{HMACSHA256(header + payload, JWT\_SECRET)}. Proves the token was issued by your server and hasn't been tampered with. \\
\bottomrule
\end{tabularx}
\end{center}

\begin{warnbox}[WARNING]
The payload is base64-encoded, not encrypted -- anyone with the token can
decode and read it (try pasting one into jwt.io). Never put passwords,
secrets, or sensitive personal data in a JWT payload. Only put identifiers
like a user ID.
\end{warnbox}

\subsection{Expiration and Signing}

Every token should carry an \texttt{exp} (expiration) claim so a stolen or
leaked token becomes useless after a fixed window. The signature is produced
using a secret key -- \texttt{process.env.JWT\_SECRET} -- that only your
server knows. If that secret leaks, anyone can forge valid tokens for any
user, so treat it like a password and never commit it to source control.

\section{The Authentication Flow}

\begin{diagrambox}[JWT Authentication Flow]
\begin{verbatim}
 1. Browser -> POST /api/auth/login (email, password)
 2. Express -> Credentials verified (bcrypt.compare)
 3. Express -> JWT generated (jwt.sign, signed with JWT_SECRET)
 4. Express -> JWT sent to browser (in an HttpOnly cookie)
 5. Browser -> stores cookie automatically, sends it on every request
 6. Browser -> GET /api/tasks  (cookie attached automatically)
 7. Middleware -> reads token, verifies with jwt.verify
 8. Middleware -> decoded payload -> req.user = user document
 9. Protected Controller -> runs, trusting req.user
\end{verbatim}
\end{diagrambox}

\section{Generating a Token}

\begin{lstlisting}[language=JavaScript]
// utils/generateToken.js
import jwt from "jsonwebtoken";

export const generateToken = (userId) => {
  return jwt.sign(
    { id: userId },
    process.env.JWT_SECRET,
    { expiresIn: "7d" }
  );
};
\end{lstlisting}

Used right after credentials are verified in the login controller:

\begin{lstlisting}[language=JavaScript]
// controllers/authController.js
import { generateToken } from "../utils/generateToken.js";

export const loginUser = async (req, res, next) => {
  try {
    const { email, password } = req.body;
    const user = await User.findOne({ email });
    if (!user || !(await bcrypt.compare(password, user.password))) {
      return res.status(401).json({ success: false, message: "Invalid credentials" });
    }

    const token = generateToken(user._id);

    res.cookie("token", token, {
      httpOnly: true,
      maxAge: 7 * 24 * 60 * 60 * 1000,
    });

    res.status(200).json({
      success: true,
      data: { id: user._id, name: user.name, role: user.role },
    });
  } catch (err) {
    next(err);
  }
};
\end{lstlisting}

\begin{notebox}[NOTE]
Full cookie configuration (\texttt{secure}, \texttt{sameSite}, and why
\texttt{httpOnly} matters) is covered in the next chapter. For now, focus on
how the token itself is created and verified.
\end{notebox}

\section{Verifying a Token}

\begin{lstlisting}[language=JavaScript]
// middleware/authMiddleware.js
import jwt from "jsonwebtoken";

export const protect = async (req, res, next) => {
  const token = req.cookies.token;
  if (!token) {
    return res.status(401).json({ success: false, message: "Not authenticated" });
  }

  try {
    const decoded = jwt.verify(token, process.env.JWT_SECRET);
    req.userId = decoded.id; // full lookup covered in Chapter 20
    next();
  } catch (err) {
    return res.status(401).json({ success: false, message: "Invalid or expired token" });
  }
};
\end{lstlisting}

\texttt{jwt.verify} checks the signature against \texttt{JWT\_SECRET} and
automatically rejects the token if the \texttt{exp} claim has passed --
either case throws, which is why it must be wrapped in \texttt{try/catch}.

\begin{tipbox}[Advanced: Access Tokens vs Refresh Tokens]
Larger production systems often split this into two tokens: a short-lived
\textbf{access token} (minutes) used on every request, and a long-lived
\textbf{refresh token} (days/weeks) used only to silently obtain a new
access token when the old one expires. This limits the damage window if an
access token leaks, at the cost of extra complexity (a refresh endpoint,
refresh token storage/rotation, revocation logic). This book uses a single
7-day token throughout to keep the implementation focused -- but know that
the two-token pattern is the natural next step for a larger production
system.
\end{tipbox}

\whatwhyhow
{A JWT is a signed token containing a user identifier and expiration, generated on login and verified on every protected request.}
{It lets the server confirm "who is this request from" without a database lookup on every request or server-side session storage, while still expiring automatically.}
{Sign with \texttt{jwt.sign(\{id\}, JWT\_SECRET, \{expiresIn\})} at login, store it in a cookie, and verify with \texttt{jwt.verify(token, JWT\_SECRET)} inside middleware before any protected controller runs.}
